{\color{maroon}
\chapter{Technical Documentation and Supporting Materials}

\section{System Usability Scale (SUS) Questionnaire}

The following questionnaire was administered to all participants following their interaction with the Consol system. Responses were collected on a five-point Likert scale ranging from "Strongly Disagree" (1) to "Strongly Agree" (5).

\subsection{SUS Questions}
\begin{enumerate}
    \item I think that I would like to use this system frequently.
    \item I found the system unnecessarily complex.
    \item I thought the system was easy to use.
    \item I think that I would need the support of a technical person to be able to use this system.
    \item I found the various functions in this system were well integrated.
    \item I thought there was too much inconsistency in this system.
    \item I would imagine that most people would learn to use this system very quickly.
    \item I found the system very cumbersome to use.
    \item I felt very confident using the system.
    \item I needed to learn a lot of things before I could get going with this system.
\end{enumerate}

\section{Post-Session Feedback Questions}

Additional qualitative feedback was collected through semi-structured interviews addressing the following areas:

\subsection{User Experience Assessment}
\begin{enumerate}
    \item How would you describe your overall experience with the semantic similarity-based feedback compared to traditional exact-match systems?
    \item Did you feel the system fairly evaluated your responses when you paraphrased or restructured information?
    \item What aspects of the progress tracking and star rating system motivated or discouraged your continued use?
    \item How intuitive did you find the note creation and practice session workflow?
    \item What improvements would you suggest for the user interface or system functionality?
\end{enumerate}



\section{Individual Participant Materials and Recollections}

This section presents the exact study materials and participant recollections used in the research validation, preserved in their original character-sensitive formatting. Each participant engaged with five core topics over a three-day testing protocol.

\subsection{Study Topics (Reference Materials)}

{\small
\textbf{TOPIC 1: CELLS}

A \textbf{cell} is the smallest structural and functional unit of life. It contains cytoplasm and a membrane, and most cells have a nucleus and organelles. Cells carry out essential functions such as energy production, reproduction, and communication. Cell biology helps us understand organisms, disease, and biotechnology.

\textbf{What a Cell Is:}\\
Examples include bacteria (E. coli), human skin cells, neurons, and unicellular protists like paramecium.\\
\textbf{What a Cell Is Not:}\\
A virus (not alive), DNA alone, a mitochondrion (only an organelle), or nonliving materials like crystals.

\textbf{Characteristics of All Cells:}\\
• Surrounded by a plasma membrane\\
• Contain ribosomes for protein synthesis\\
• Possess DNA as genetic material\\
• Capable of homeostasis\\
• Can divide to reproduce\\
• Respond to stimuli

\textbf{Two Main Types of Cells:}

1. \textbf{Prokaryotic Cells} – No true nucleus, no membrane-bound organelles, circular DNA, smaller size; divide by binary fission (bacteria, archaea).

2. \textbf{Eukaryotic Cells} – Have a nucleus, many organelles, linear DNA; divide by mitosis/meiosis (animals, plants, fungi, protists).

\textbf{Basic Cell Structures:}\\
• Plasma Membrane – Controls movement of substances.\\
• Cytoplasm – Gel-like substance filling the cell.\\
• Nucleus – Control center containing DNA.\\
• Mitochondria – Powerhouses producing energy.\\
• Endoplasmic Reticulum – Transport system.\\
• Golgi Apparatus – Processing and packaging center.\\
• Ribosomes – Protein synthesis sites.\\
• Flagella/Cilia – Movement.

\textbf{Cell Theory:}\\
Developed by Hooke, Leeuwenhoek, Schleiden, Schwann, and Virchow.\\
Three principles:

1. All living things are made of cells.

2. The cell is the basic unit of life.

3. All cells come from pre-existing cells.

---

\textbf{TOPIC 2: THE ENDOCRINE SYSTEM}

The endocrine system is a network of glands that produce \textbf{hormones}, the chemical messengers that allow cells and organs to communicate. These hormones help control growth, metabolism, reproduction, stress response, and mood.

A \textbf{gland} is an organ that makes and releases hormones into the bloodstream. Together, endocrine glands coordinate major body functions by producing hormones, regulating how much is released, and sending them into the blood so they reach their target organs.

The endocrine system includes many glands located throughout the body. In the brain are the \textbf{hypothalamus}, which links the endocrine and nervous systems, and the \textbf{pituitary gland}, the "master gland" that controls other glands. The \textbf{pineal gland} makes melatonin for sleep.

In the neck, the \textbf{thyroid} controls growth and metabolism, while the \textbf{parathyroid glands} regulate calcium and phosphorus for bone health. The \textbf{thymus} helps develop immune cells in childhood. The \textbf{adrenal glands}, on top of the kidneys, produce adrenaline and corticosteroids.

Some organs also function as endocrine glands. The \textbf{pancreas} makes insulin and glucagon to regulate blood sugar. In females, the \textbf{ovaries} produce estrogen and progesterone. In males, the \textbf{testes} make testosterone.

---

\textbf{TOPIC 3: APPLIED SCIENCE}

Applied science is the use of scientific knowledge to solve real problems. It involves analysis, experimentation, and improving solutions through repeated testing. Applied science focuses on specific problems, uses existing scientific principles, and aims to create practical and useful results.

A key feature of applied science is its \textbf{problem-oriented focus}. Scientists begin with a clear need—such as improving internet speed, creating safer medicines, or designing more efficient machines.

Applied science also involves \textbf{experimentation and refinement}. Early versions of solutions are rarely perfect. Engineers and scientists test prototypes, gather data, and improve the design until it is efficient, reliable, and cost-effective.

Applied science appears in many fields:\\
• \textbf{Computer science} – creating software, algorithms, secure networks, and user-friendly apps.\\
• \textbf{Biomedical engineering} – developing medical devices, prosthetics, and diagnostic tools.\\
• \textbf{Electrical engineering} – designing circuits, power systems, and electronic devices.

---

\textbf{TOPIC 4: PSORIASIS}

Psoriasis is a long-term skin condition that causes itchy, scaly patches that commonly appear on the knees, elbows, scalp, and trunk. It can be painful, affect sleep, and make it hard to concentrate. The disease usually comes in cycles, with flare-ups lasting weeks or months before calming down again.

There are several types of psoriasis. \textbf{Plaque psoriasis}, the most common, causes dry, raised, scaly patches. \textbf{Nail psoriasis} affects fingernails or toenails, causing discoloration, pitting, abnormal growth, or nail separation.

Psoriasis is believed to be an immune system disorder where infection-fighting cells attack healthy skin by mistake, causing rapid skin cell growth. Common triggers include:\\
• Infections (especially strep throat)\\
• Weather changes\\
• Smoking\\
• Alcohol consumption\\
• Certain medications\\
• Sudden withdrawal of steroids

---

\textbf{TOPIC 5: EARTHQUAKE PREPAREDNESS}

\textbf{Definition:}\\
An earthquake is the shaking of the Earth's surface caused by sudden movement in the Earth's crust. These movements release energy in the form of seismic waves, which travel through the ground and can be detected by instruments called seismographs.

\textbf{Causes:}

• \textbf{Tectonic Plate Movements} – Most earthquakes occur along plate boundaries due to collisions, separations, or sliding past one another.\\
• \textbf{Volcanic Activity} – Magma movement beneath the surface can create small earthquakes.\\
• \textbf{Human Activity} – Mining, reservoir construction, and underground explosions can sometimes trigger quakes.

\textbf{Effects:}

• \textbf{Physical Damage:} Collapse of buildings, bridges, roads, and other infrastructure.\\
• \textbf{Human Impact:} Injuries, loss of life, and displacement of communities.\\
• \textbf{Environmental Impact:} Landslides, tsunamis, and changes in the landscape.
}

\subsection{Individual Participant Recollections}

The following sections present each participant's exact recollections as written during the testing sessions, preserving original formatting, spacing, and character layout.

\subsubsection{Participant 1 Recollections}

{\footnotesize
\textbf{TOPIC 1: THE CELL  day 1}


A cell is the smallest structural and functional unit of life. it also contains energy, communication. the cell is very important in the human history because of bong bong marcos and his minions grr😌😌😌


TYPES OF CELLS:🧫
•Brain cells
•Blood cells
•Plant cells
•Animal cells
•White cells
•Sperm cells

\textbf{TOPIC 2: ENDOCRINE day 1}
 
The endocrine system is a network of glands that produce hormones, the chemical messengers that allows cells and organs to communicate. These hormones help control growth, metabolism, reproduction, stress response, and mood. If the endocrine system is unhealthy, a person may struggle with puberty, fertility, energy levels, weight control, or blood sugar balance.


A gland is a organ that makes and releases hormones into bloodstream. together endocrine glands

\textbf{TOPIC 3: day 1}


applied science is the use of scientific knowledge to solve real problems. it is not just using science automatically instead, it involves analysis experimentation and improving solutions through repeated testing applied science focuses on specific problems uses existing scientific principles and aims to create practical and useful results 


A key feature of applied science is its problem oriented focus scientists begin with a clear need such as improving internet speed creating safer medicines or designing more efficient machines instead of discovering new

\textbf{TOPIC 4: PSORIASIS day 1}


psoriasis is a long-term skin condition that causes itchy scaly on the knees elbows scalp and trunk it can be painful affect sleep and make it hard to concentrate the disease usually comes in cycle with flare ups lasting weeks or months before calming down again psoriasis is not contagious but it can greatly affect comfort and daily life symptoms vary from person to person many develop a patchy rash that may look like different depending on skin color pink or red patches with grey scales on white

\textbf{TOPIC 5: day 1}


An earthquake is the shaking of the earthz surface caused by sudden movement in the earths crust these movements release energy in the form of seismic waves which travel through the ground and can be detected by instruments called seismographs


CAUSES:
•tectonic plate movements 
•volcanic activity 
•human activity 


EFFECTS:
•physical damage 
•human impact 
•environmental impact 


SAFETY MEASURES:
•build earthquake resistant structures 
•prepare emergency kits

\textbf{TOPIC 1: day 2}


How cells move:
cells move using flagella whip like tails cilia short beating hairs or amoeboid movement using pseudopodia
 


How cells reproduce:
binary fission prokaryotes 
mitosis makes identical somatic cells
meiosis produces gametes with half the chromosome number


what cells do:
cells carry out life processes 
metabolism energy reproduction molecule transport growth repair communication reproduction and defense (immune cells)


orgin of cells
life likely began throu abiogenesis with self re

\textbf{TOPIC 2:}


in the neck the thyroid controls growth and metabolism whike the parathyroid glands regulate calcium and phosphorus for bone health the thymus located between the lungs helps develop immune cells in childhood The adrenal glands on top of the kidneys produce adrenaline and corticosteroids which affect stress response and metabolism 


some organs also function as endocrine glands The pancreas makes insulin and glucagon to regulate blood sugar in females the ovaries produce estrogen and progesterone for puberty the menstrual cycle and pregnancy.

\textbf{TOPIC 3:}


All applied science follows a structured process


•COMPUTER SCIENCE - creating software algorithms secure networks and user friendly apps
•ELECTRICAL ENGINEERING - building circuits electronics and communication systems like 5G
•MECHANICAL ENGINEERING - designing

\textbf{TOPIC 4:}


These are the types of PSORIASIS 


Plaque psoriasis 
Nail psoriasis 
guttate psoriasis 
inverse psoriasis 
pustular 
erythrodermic


PSORIASIS is believed to be an immune system disorder where infection systems disorder where infection fighting celks attack healthy skin by mistake causing rapid skin cell growth genetic and environmental factors both contribute common triggers include infection

\textbf{TOPIC 5:}


CONCLUSIONS:
Earthquake are natural hazards that can have devastating consequences understanding their causes and effects and being prepared are essential to reducing risk and saving lives 


PRECAUTIONARY MEASURES FOR EARTHQUAKES


Before an Earthquake (PREPAREDNESS)


secure yor home
emergency kit
safety plan
education and drills

\textbf{TOPIC 1, day 3}


Cells are the smallest form of life or like the most basic unit of living being. Cells come from many forms. Cells can make up all kinds of tissues, or organelles. A cell can be for humans, for animals, and for plants. Prokayotes are cells that are in bacteria, and eukaryotes are cells are in animals and plants. There are many parts of a cells: like the powerhouse of the cell which is the mitochondria, and the golgi apparatus, the plasma membrane, the cell wall, etc

\textbf{TOPIC 2, day 3}


The endocrine system is part of the many systems in the human body. This system is for hormones and various chemicals in benefit for the body and its functions. The glands in part of this system include: the thyroid, the parathyroid, the pancreas, etc. These help in things such as keeping bone health good, and develop immunity for diseases. 

Endocrine is one of the hormones. Disorders can develop when there are excess hormones in the body. 

Pancreas makes glucagon to regulate blood sugar and in females, the varies produces estrogen and progesterone for the menstruation.

\textbf{TOPIC 3, day 3}




Applied Science is deals with real world problems. Science has a lot of areas but applied science deals with fields that have practical use. Computers, electric vehicles, and all the way to chemicals are parts of applied sciences. 

To name a few: 
Electrical engineering is for building circuitry and dels with voltages in electricity
Computer Engineering is for making hardware and combining it with digital chu
Material science is for making objects with specific intended uses like ceramics and plastics
Data science is for analyzing information that is scattered and making it make sense in our eyes.

Experimenting is crucial to sciences, but especially in applied science, wherein it needs to know how the result will become.

\textbf{TOPIC 4, day 3}




 One of the skin diseases  that cause patchy scaly skin is psoriasis, and it comes from bacteria or fighting of the immune system that becomes a disorder. 


The symptoms include infections like sore throat, dry skin, burns, and many other.
Causes from cold or dry weather in non humid places or unhealthy habits or even stopping to use steroid medication. It is genetic as it can affect anyone.

\textbf{TOPIC 5, day 3}










During an earthquake:
Duck, Cover, and Hold. Duck toget under cover, and hold onto stable objects as the ground shakes. 

WHy and earthquake happens
It is mostly because of plate movement that we cannot see, and the shock becomes a cause or the epicenter of the tremors. 

Volcanic activity is also one of the reasons, because eruptions are coming from magma inside the volcano going to burst, making excessive kinetic energy. 

Humans can also trigger the ground to collapse and shift accidentally. Sometimes man made activities like digging and mining.


Before an earthquake we must gather enough materials in emergency bags so access to it is quick enough in evacuation. Anchoring or nailing things to the wall securely is also good for avoiding collapse inside homes that might trap occupants.
}

\subsubsection{Participant 2 Recollections}

{\footnotesize
\textbf{TOPIC 1, DAY 1}


The Cell 
* Is the basic unit of living organisms
* Carries out functions of the body


Cell Biology - study of cells


Characteristics
* Presence of cytoplasm
* Presence of ribosomes
* Undergo cell division
* Reacts to stimuli
* Contains genetic material


Types:
Prokaryotic Cell
* No nucleus
* Lacking organelles
* Present in bacteria
Eukaryotic Cell
* Has nucleus
* Has organelles
* Present in animals, plants, fungi


Cell Structures
1. Cytoplasm
2. Cell Membrane
3. Mitochondria
4. Endoplasmic reticulum
5. Golgi apparatus
6. Ribosome


Cell Theory - explains the basic principles of cells
3 Principles:
1. Cells are the basic unit of life
2. Life is made of one or more cells
3. Cells arise from pre-existing cells

\textbf{TOPIC 2, DAY 1}


Endocrine System - makes and manages hormones
Hormones - chemicals the body releases for growth and development
Gland - organ that produces certain hormone


Glands:
1. Pituitary - master gland
2. Adrenal - produces adrenaline
3. Thymus - controls pituitary gland
4. Thyroid - 
5. Parathyroid

\textbf{TOPIC 3, DAY 1}


Applied Science 
* process of experimenting, analyzing, and making conclusions from data 


* Applied science adds on to what pure science achieves


Real World Applications:
1. Biomedical Engineering
2. Computer Science
3.

\textbf{TOPIC 4, DAY 1}


Psoriasis
* Thought to be an immune disease that abnormally speeds up the rate of production of skin cells
* Symptoms include: rashes, patchy skin


Triggers:
1. Weather
2. Smoking
3. Alcohol consumption
4. Certain medication
5. Wounds on the skin

\textbf{TOPIC 5, DAY 1}


Earthquakes - seismic activity that causes shaking and breaking of earth's surface


Causes:
1. Plate Tectonics - a natural phenomenon caused by convection currents in the mantle, making tectonics plates to move
2. Volcanic Eruption - magmatic rocks exit out of natural vents 
3. Man-made Activities - eg. mining

\textbf{TOPIC 1, DAY 2}


The Cell
* Basic unit of life




Characteristics
* Presence of cytoplasm
* "              " of cell membrane
* Ribosome
* Genetic material
* Reacts to stimuli
* Undergo cell division
* Capable of homeostasis


Cell Structure
1. Mitochondria
2. Cytoplasm
3. Cell membrane
4. Ribosome
5. Golgi apparatus
6. Endoplasmic reticulum


Cell Theory - explains the principles of cells
Principles: 
1. All life is made of one or more cells
2. The cell is the basic unit of life
3. Cells arise from pre-existing cells


Cell Reproduction
1. Binary Fission
2. Mitosis
3. Meiosis


Cell Mechanisms

\textbf{TOPIC 2, DAY 2}


Endocrine System - consist of glands that make and manage hormones
Hormones - regulate growth and development
Gland - organ that secretes hormones


Glands:
Hypothalamus - manages pituitary gland
Pituitary Gland - master gland
Thyroid - metabolism
Parathyroid - manages bone health
Adrenal - adrenaline
Pineal - melatonin


Organs:
Pancreas
Testes

\textbf{TOPIC 3, DAY 2}


Applied Science - process of analysis, experimentation, and refinement of findings


Characteristics: 
* Problem oriented
* Application of existing knowledge
* Practical and Efficient


Real-life Applications
1. Computer Science
2. Biomedical engineering
3. Electrical engineering
4. Chemical engineering
5. Aerospace engineering
6. Robotics
7. Mechanical engineering


Challenges:
* Complexity of real world problems
* Resource constraint
* Ethical considerations
* Technological change

\textbf{TOPIC 4, DAY 2}


Psoriasis 
 - skin disease that causes itchiness, rashes, and patchy skin.
- has no cure
- skin cells grow faster than usual


Triggers
1. Infection
2. Weather
3. Skin injury
4. Smoking
5. Alcohol consumption
6. Certain medication


See a doctor if:
* Becomes widespread
* Causes pain
* Doesnt improve with treatment


Complications:
* Psoriatic arthritis
* Eye conditions
* Obesity
* Type 2 diabetes
* High blood pressure 
* Cardiovascular disease
* Autoimmune diseases

\textbf{TOPIC 5, DAY 2}


Earthquakes - shaking of earth's surface due to movement of earth's crust


Causes:
1. Tectonic plate movement - occurs along plate boundaries
2. Volcanic eruption - movement of magma underneath earth
3. Human activity - eg. mining, underground explosions, and construction


Effects:
Physical Damage - Infrastructure damage
Human impact - Loss of life
Environmental impact - Landslide, Tsunamis

\textbf{TOPIC 1, DAY 3}


The Cell
* Is the basic fundamental unit of life responsible for carrying out multiple functions in the body
It is:
* A bacterium
* A skin cell
* A neuron
It is not:
* A virus
* DNA or RNA alone


Characteristics
* has plasma membrane
* Has cytoplasm
* Has ribosomes for protein synthesis
* Possesses DNA or RNA
* Capable of homeostasis
* Can divide 
* Responds to stimuli


Types:
1. Prokaryotic Cells
* No nucleus
* Lacking organelles
* Found in bacteria
2. Eukaryotic Cells
* Has nucleus
* Has organelles
* Found in animals, plants, fungi 


Cell Structures
Plasma membrane
Cytoplasm
Golgi apparatus
Endoplasmic reticulum

\textbf{TOPIC 2, DAY 3}


Endocrine System - network of glands that make and manage hormones
Hormones - chemicals that affect growth and development
Glands - organs that secrete hormones


Glands:
1. Hypothalamus - controls pituitary gland
2. Pituitary Gland - the master gland
3. Pineal Gland - melatonin for sleep
4. Thyroid Gland - metabolism
5. Parathyroid - bone health
6. Thymus - immunity
7. Adrenals - adrenaline 


Organs: 
1. Pancreas
2. Ovaries (for women)
3. Testes (for men)

\textbf{TOPIC 3, DAY 3}


Applied Science - process of analysis, experimentation, and refinement of discoveries


Characteristics:
Problem-oriented
Application of Existing Knowledge
Experimental


Applications:
Electrical engineering
Mechanical engineering
Chemical engineering
Software engineering
Biomedical engineering


Structured approach of Applied Science
Problem definition
Literature review
Prototyping
Testing
Analysis

\textbf{TOPIC 4, DAY 3}


Psoriasis - skin disease that causes rash, itchiness, and scaly patches
* Long term, has no cure, and can be painful


Symptoms: 
* Patchy rash
* Scaling spots
* Cracked skin
* Burning sensation of skin


Types: 
1. Plaque
2. Nail


See a doctor if:
* Becomes widespread
* Causes pain
* Doesnt improve with treatment


Triggers
* Infection
* Weather
* Smoking
* Alcohol consumption
* Wounds on skin
* Certain medication

\textbf{TOPIC 5, DAY 3}


Earthquakes 
* Sudden shaking of earth's surface, releasing seismic waves that can be detected by seismographs


Causes: 
Plate Tectonics - earthquakes occur along plate boundaries or faults
Volcanic Eruption - movement of magma beneath the earth
Human activity - eg. mining and construction can cause disturbances underneath


Effects:
* Physical damage - damage to infrastructure
* Human impact - casualties and injuries
* Environmental impact - leads to landslides and tsunamis

\textbf{TOPIC 1, DAY 4}
The Cell
* Is the fundamental unit of life, it is the foundation of organisms


It is: 
* Bacterium
* A neuron
* A skin cell


It is not: 
* A mitochondria 
* A virus
* Genetic material alone


Characteristics
All cells have:
* Plasma membrane
* Cytoplasm
* Ribosomes
* DNA/RNA 
* Capable of homeostasis


Types: 
1. Prokaryotic 
* No nucleus
* Lacking organelles
* Found in bacteria
2. Eukaryotic 
* Has nucleus 
* Various organelles
* Found in animals, plants, fungi


Cell mechanisms:
* Flagella - whip-like tails
* Cillia - hair-like structures
* Amoeboid movement


Cell reproduction:
* Binary fission (for prokaryotes)
* Meiosis
* Mitosis

\textbf{TOPIC 2, DAY 4}
 
Endocrine System 
* Network of glands that make and manage hormones in the body
Gland - organ that secretes hormone
Hormone - chemical that aids growth and development


Endocrine Glands:
Hypothalamus - manages pituitary 
Pituitary gland - master gland
Pineal gland - produces melatonin to signal sleep
Thyroid gland - metabolism
Thymus - immunity


Endocrine organs:
* Pancreas
* Ovaries (in women)
* Testes (in men)


Disorders: 
Acromegaly 
Addison's disease
Hyperthyroidism

\textbf{TOPIC 3, DAY 4}


Applied Sciences
* Application of scientific discoveries, involving analysis, experimentation, and refining.


Characteristics: 
* Problem oriented
* Application of existing knowledge
* Creating solutions
* Experimental


Domains:
1. Chemical Engineering
2. Softwaire Engineering
3. Electrical Engineering
4. Aerospace Engineering
5. Data analytics


Methodologies:
1. Defining the problem
2. Literature review
3. Prototyping
4. Testing
5. Opitimizing

\textbf{TOPIC 4, DAY 4}


Psoriasis 
* Immune system problem that causes skin cells to grow faster than usual


Triggers: 
* Infection
* Weather
* Smoking
* Excessive alcohol consumption
* Certain medication


See a doctor if:
* Widespread
* Painful
* Doesnt improve with treatment


Psoriasis may lead to:
* Obesity
* Type 2 diabetes
* High blood pressure

\textbf{TOPIC 5, DAY 4}


Earthquakes - sudden shaking of the earth's surface due to movement of earth's crust


Causes:
1. Plate tectonics - earthquakes usually occur at plate boundaries or faults
2. Volcanic activity - movement of magma underneath earth
3. Human activity - eg. mining and construction that disrupt earths interior


Effects:
* Loss of life
* Injuries
* Infrastructure damage
* Other natural disasters eg. tsunami, landslide, volcano eruption

\textbf{TOPIC 1, DAY 5}


The Cell
* Is the basic fundamental unit of life responsible for carrying out multiple functions in the body
It is:
* A bacterium
* A skin cell
* A neuron
It is not:
* A virus
* DNA or RNA alone
* A mitochondrion
Characteristics
* has plasma membrane
* Has cytoplasm
* Has ribosomes for protein synthesis
* Possesses DNA or RNA
* Capable of homeostasis
* Can divide 
* Responds to stimuli
Types:
1. Prokaryotic Cells
* No nucleus
* Lacking organelles
* Found in bacteria
2. Eukaryotic Cells
* Has nucleus
* Has organelles
* Found in animals, plants, fungi 
Cell Structures
1. Plasma membrane
2. Cytoplasm
3. Golgi apparatus
4. Endoplasmic reticulum
5. Ribosome
6. Mitochondria


Cell Mechanisms:
* Flagella - whip-like tails
* Cillia - hair-like structure
* Amoeboid movement


Cell theory - describes the basic properties roles of all cells

\textbf{TOPIC 2, DAY 5}


Endocrine System 
* Network of glands that make and manage hormones in the body
Gland - organ that secretes hormone
Hormone - chemical that aids growth and development


Endocrine Glands:
Hypothalamus - manages pituitary 
Pituitary gland - master gland
Pineal gland - produces melatonin to signal sleep
Thyroid gland - metabolism
Thymus - immunity
Parathyroid - bone health/ calcium & phosphorus levels
Adrenals - produces adrenaline, responsible for fight or flight


Endocrine organs:
* Pancreas
* Ovaries (in women) - estrogen, progesterone
* Testes (in men) - testosterone


Disorders: 
Acromegaly 
Addison's disease
Hyperthyroidism 
Hypopituitarism 
Osteoperosis

\textbf{TOPIC 3, DAY 5}
Applied Science - process of analysis, experimentation, and refinement of discoveries


Characteristics:
* Problem-oriented
* Application of Existing Knowledge
* Experimental
* Practical & Efficient


Applications:
* Electrical engineering
* Mechanical engineering
* Chemical engineering
* Software engineering
* Biomedical engineering
* Aerospace engineering
* Data analytics


Structured approach of Applied Science
Problem definition
Literature review
Prototyping
Testing
Analysis
Refinement
Validation

\textbf{TOPIC 4, DAY 5}


Psoriasis - skin disease that causes the skin cells to grow faster than normal


Symptoms: 
* Patchy rash
* Scaling spots
* Cracked skin
* Burning sensation


Types:
* Nail psoriasis
* Guttate psoriasis
* Pustular psoriasis
* See a doctor if:
* Widespread across body
* Painful
* Gets worse despite treatment
Triggers:
* Infection
* Dry weather
* Skin injuries
* Smoking
* Certain medication
* Excessive alcohol consumption 


This may lead to:
* Cardiovascular disease
* High blood pressure
* Type 2 diabetes

\textbf{TOPIC 5, DAY 5}


Earthquakes 
* Natural event wherein sudden shaking of earth's surface, caused by movement of earth's crust 
*  releasing seismic waves that can be detected by seismographs


Causes: 
Plate Tectonics - earthquakes occur along plate boundaries or faults, 
Volcanic Eruption - movement of hot magmatic material beneath the earth
Human activity - eg. mining and construction can cause disturbances underneath


Effects:
* Physical damage - damage to infrastructure
* Human impact - casualties and injuries
* Environmental impact - leads to landslides and tsunamis


Things to prepare:
* Evacuation plan
* Emergency kit
}

\subsubsection{Participant 3 Recollections}

{\footnotesize
\textbf{TOPIC 1: The Cell | Day 1}


A cell is the smallest structural and functional unit of life.


What Cell is:
-Examples include; neurons, unicellular, paramecium, human skin and bacteria. 


What Cell is Not: 
-DNA, virus, mitochondrion.


Two main type of cells:
-Prokaryotic Cells
-Eukaryotic Cells


What cells do:
Cells carry out life process, like; metabolism, energy production,

\textbf{TOPIC 2: The Endocrine System | Day 1}


It is a network of glands that produces hormones, the chemicals that allow cells and organs to communicate with each other.

\textbf{TOPIC 3: Applied Science | Day 1}


Applied Science is the use of scientific knowledge to solve real problems. 


Applied Science focuses on specific problems, uses existing scientific principles and aims to create practical and useful results.


Applied Science also involves experimentation and refinement. 


Applied Science appears in many fields: 
•Computer Science 
•Electrical engineering 
•Mechanical engineering 
•Mater

\textbf{TOPIC 4: Psoriasis | Day 1}


Psoriasis is a long-term condition that causes itchy, scaly patches that commonly appear on the knees, elbows, scalp, and trunk. 


There are several types of psoriasis:
•Plaque 
•Psoriasis 
•Nail Psoriasis 
•Guttate Psoriasis 
•Inverse Psoriasis 
•Pustular Psoriasis 
•Erythrodermic Psoriasis

\textbf{TOPIC 5: Earthquake Preparedness | Day 1}


An earthquake is the shaking of the earth's surface caused by sudden movement in the earth's crust. 


Causes:
•Tectonic Plate Movements 
•Volcanic Activity 
•Human Activity 


Effects: 
•Physical Damage 
•Human Impact 
•Environmental Impact

\textbf{TOPIC 1 | DAY 2}


 A cell is the basic unit of life. It is the smallest structure that can perform all the processes needed for a living organism to survive.


All living things are made up of cells.
Some organisms have one cell (unicellular), like bacteria.
Others have many cells (multicellular), like humans, animals, and plants.


Cells carry out important functions like:
•Breathing
•Growing
•Getting energy
•Responding to the environment


Examples:
•Muscle cell (in humans)
•Nerve cell
•Blood cell
•Plant leaf cell
•Bacterial cell

\textbf{TOPIC 2 | DAY 2}


The endocrine system includes many glands located throughout the body. In the brain are the hypothalamus, which links the endocrine and nervous systems, and the pituitary gland, the "master g

\textbf{TOPIC 3 | DAY 2}


Applied Science is the bridge between scientific theory and real-world innovation — turning knowledge into technology that improves everyday life. 


All applied science follows a structured process: define the problem, study existing knowledge, design a solution, test it, improve it, and confirm that it works in real condition.

\textbf{TOPIC 4 | DAY 2}


Psoriasis is believed to be an immune system disorder where infection-fighting cells attack healthy skin by mistake, causing rapid skin cell growth. Genetics and environmental factors both contribute. 


Common Triggers Include:
•Infections (especially strep throat)
•Cold, dry weather 
•Cuts, scrapes or sunburn 
•Smoking or second hand smoke 
•Heavy alcohol use 
•Certain Medications 
•Certain withdrawal of steroids

\textbf{TOPIC 5 | DAY 2}


Earthquakes are natural hazards that can have devastating  consequences. Understanding their causes and effects and being prepared are essential to reducing risk and saving lives.

\textbf{TOPIC 1 | DAY 3}








Cells are the smallest structural unit of life and it makes up tissues, then it makes up organs, and then organ systems. Multicellularity and specialization of cells allow it to form other kinds of organs or tissues for specific purposes in the human body. Cells mostly have the mitochondria which is the powerhouse of the cell, and not all have nucleus. All cells come from pre-existing cells, and they contain DNA, which makes us inherit things from our predecessors. Cell biology is what deals with this area of science. 

Cells carry out important functions like:
•Growing
•Getting energy
•Reproduction










Examples:
•Muscle cell
•Plant leaf cell
•Red Blood cell
•White Blood cells
•Bacterial cell
•Nerve cell

\textbf{TOPIC 2 | DAY 3}


The endocrine system include many glands that make hormones for growth and energy, and other involuntary functions in the body. 

Examplse: 

pituitary gland is the "master gland" because it controls other glands too. 
The thyroid controls how your body uses energy
The parathyroid is for calcium
Adrenal gland is near the kidneys and gives us adrenaline which makes us react to stimuli more intensely
Pancreas helps control blood sugar with insulin
Ovaries and testes make hormones for women and men ( estrogen and progesterone for women, and testosterong for men) (can sometimes lead to disorders like PCOS) 
All endocrine disorders can be influenced by a range of reasons from genetics to infections.

\textbf{TOPIC 3 | DAY 3}


 Applied science is about using scientific knowledge to solve real world problems using analyzing and experimenting while also giving new solutions by continuously testing and doing the scientific method. These work because of the idea of fixing problems that exist in the world already, so that these are addressed especially when they in demand of a solution. Applied science works more on things that are already available in the present, so protocols are set to have seamless transitions.


Fields of Applied science:
* Computer science
* Engineering
* Electrical engineering
* Mechanical engineering
* Biomedical engineering
* Data science

\textbf{TOPIC 4 | DAY 3}










Psoriasis is thought to happen because the immune system gets confused and starts attacking healthy skin cells by accident. This makes the skin cells grow too fast. Both a person's genes and things in the environment can take part in causing it.
Common things that trigger it are:
strep throat
dry weather
sunburn
smoking/smoke
Alcohol

Inverse psoriasis affects skin where it folds. Plaque psoriasis affects loose skin like in elbows and knees where it raises a little bit. Nail psoriasis affects fingers and toes. Erythrodermic psoriasis gives burns or itches when it causes peeling.

\textbf{TOPIC 5 | DAY 3}

Possible Causes of earthquake

• Volcanic activity – magma shifting under the Earth
• Human actions – mining, blasting underground, big constructions
• Tectonic plates moving, mostly where the plates meet


What to do in an earthquake
• Duck cover and hold
• make sure to stay calm
• in streets or cities, do not go outside because of falling debris from buildings
• you can move to very open areas with little to no tall collapsing things

After the earthquake:
• check for injuries and do first aid
• watch out for gas leaks or open live wires
}

\subsubsection{Participant 4 Recollections}

{\footnotesize
\textbf{TOPIC 1 DAY 1}
A cell is a smallest structural and functional unit in life, example of cell a bacterial cell e iol. Some cells are found in the organs or body like blood cell

\textbf{TOPIC 2 DAY 1}
The endocrine system a network that produce hormones, chemicals that message allowing cells and organs to communicate. Hormones can also help control growth, metabolism, reproductive, stress, response,and mood. Gland is also an organ

\textbf{TOPIC 3 DAY1}
Applied science used for scientific knowledge to solve problems. It also focuses on specific problems, like specific principals and to create practical and useful information 


Problem - oriented focus Science

\textbf{TOPIC 4 DAY 1}
Psoriasis a long term skin condition causing itchy,scaly patch that commonly appear  on tha knees, elbows scalp and trunk it is painful affect sleep and make it hard to concentrate 


Types of psoriasis plaque psoriasis cause is dry and scaly patch 


Nail psoriasis finger or toenail

\textbf{TOPIC 5 DAY 1}
Earthquake is the shaking of the earths surface. A sudden movement of the erths crust 


Cause 
Tectonic plate movement plate boundaries due to collisions or separation or sliding past on one another 


Volcanic activity magma movement beneath th surface 


Physical damage building collapsing and people getting injure

\textbf{TOPIC 1 DAY 2}
Characteristics of cells are surrounded by palm membrane,contain cytoplasm,have genetic material like dna rna


Two main types of cells
Prokaryotic calls non true nucleus, non membrane bound organelles and eukaryotic cells have many nucleus organs and linear dna

\textbf{TOPIC 2 DAY 2}
Endocrine system has many glands located in the body,such as Hypothalamus are located in the brain and pituitary gland that controls other glands 


Thyroid controls growth

\textbf{TOPIC 3 DAY 2}
applied science appears in many fields such as computer science creating software algorithms and friendly apps e


Electrical engineering building, electronics and communication system 


Mechanical engineering designing engines vehicle robots and machines


Data science can analyze large datas and build ai

\textbf{TOPIC 4 DAY 2}
psoriasis can be developed through family history and smoking can increase risk
It can also lead to completions like eye problem obesity high blood pressure hear disease and depression and other illnesses 


Trigger includes
Infection

\textbf{TOPIC 5 DAY 2}
earthquakes are natural shaking of the earths surface that can lead to a dangerous hazard and can have detestating consequences 


During an earthquake we should stop drop and hold on stay calm and and not panic we should 


After an earthquake we should check for injuries

\textbf{TOPIC 1 DAY 3}


A cell is the smallest unit of life and has a cytoplasm and a cell membrane. But not all have nucleus and organelles. 

Viruses, DNA, and crystals are not cells or are not amd out of cells. 

Two main types of cells:
Prokaryotes are cells with no nucleus and have a circular dna. No membrane organelles
Eukaryotes have a nucleus and linear dna. Eukaryotes can do mitosis or meiosis.



Cells move through flagella and cilia which are like tiny hairs.

\textbf{TOPIC 2 DAY 3}


The endocrine system is what produces hormones in the body. It controls growth, or metabolism, or reproduction, and many other.,

Different glands:
Hypothalamus
Thyroid
Parathyroid
Testes
Ovary
Pancreas

Risk factors of endocrine imbalance can involve age or genetics.

\textbf{TOPIC 3 DAY 3}






Applied science is using scientific knowledge to solve real problems. 
It is practical and "problem-focused" approach. It is different from other science because it is applied in real life and uses the scientific method to study existing info and test and improve. And it continues this cycle.


Examples are vaccine, electric cars, and data centers.


Types of applied scienc:
Data science chemical science
Mechanical science
Electrical engineering
Computer engineering

\textbf{TOPIC 4 DAY 3}








Psoriasis is a kind of skin disease that appears as scaly and dry patches or cracked skin. Sometimes it burns or itches. 

It is caused by either infections, or weather, or unhealthy habits, or even from genetics.

Types of psoriasis:
Pustular
Inverse
Nail
Plaque

See a doctor when symptoms show

\textbf{TOPIC 5 DAY 3}


DUck cover and hold during earhquare. Prepare emergency kits before earthquakes or at all times in case of earthquakes. Participate in earthquake drils and have evacuation plans. Be careful of falling objects and being near tall structures.


Earthquae is from mining and digging but mainly from the shaking of the earths plates due to techntonic movement.
}

\subsubsection{Participant 5 Recollections}

{\footnotesize
\textbf{TOPIC 1: THE CELL | DAY 1}


A cell is a smallest and a fundamental unit of life that can carry out all the vital 
biological functions. A cell is not a DNA, mitochondria, crystal, water droplet, and a virus.
The cell helps us to live and the cell is not only found in humans but in everything, it can be a plant, animal, bacteria, and a fungi. There are two types of cells: the Prokaryotic Cells and Eukaryotic Cells. Prokaryotic cells are found in bacterias, while the Eukaryotic Cells are found in animals, plants, and fungi. The structure of the cells has (mixed version): cell wall, golgi apparatus, mitochondria,

\textbf{TOPIC 2: THE ENDOCRINE SYSTEM | DAY 1}


The endocrine gland plays a vital role in our body that helps our body function normally.
This helps us balance and maintain our body to have a healthy life. If the endocrine gland is unhealthy it may affect our body, mood, and balance. The parts of the endocrine system are the pilinial gland, pituitary gland, thyroid,

\textbf{TOPIC 3: APPLIED SCIENCE | DAY 1}


In today's modern world we often cling to Science for more discoveries, wonders, and awe-inspiring information. Science helps us to produce many inventions and innovations. Applied Sciences appears in many fields like Computer Science, Mechanical Engineering, Biomedical Engineering, Aerospace Engineering, and Biomedical Engineering.

\textbf{TOPIC 4: PSORIASIS | DAY 1}


Psoriasis is a complex skin infection that may affect the person's well-being. And it may have a cure for it but we need to be careful because prevention is better than cure. :D And the possible ways to prevent this is to prevent sun exposure, eating healthy, and wearing sunblock, etc. Moreover, there is nothing to worry about because there are professional dermatologists that can help you through it :) And pray.

\textbf{TOPIC 5: EARTHQUAKE PREPAREDNESS | DAY 1}


Earthquakes are known as a destructive calamity that affects many living things and infrastructures. Earthquakes will happen naturally and unnaturally. And we may record and track earthquakes in a seismograph and it is monitored by the seismologist. We must be prepared and be ready for unexpected earthquakes.

\textbf{TOPIC 1: THE CELL | DAY 2}


The cell is the basic unit of life and it helps all living organisms live and experience life. All living things are made up of cells. Some organisms have only one cell (unicellular), while others have millions (multicellular). Each cell has parts with specific jobs, like the nucleus that controls activities, the cell membrane that protects the cell, and the cytoplasm where many processes happen. Cells work together to keep the body alive and functioning. There are two types of cells we have: the prokaryotic and eukaryotic cells. These cells are not always found in all living organisms, these cells have their own living organisms.


Cell stru

\textbf{TOPIC 2: THE ENDOCRINE SYSTEM | DAY 2}
The endocrine system is composed of different glands which secretes hormones 
that regulate metabolism, growth and development, mood and reproduction. If  the endocrine system malfunctions our body may have problems in reproduction, 


THE ENDOCRINE SYSTEM:
* PINEAL GLAND
* PITUITARY GLAND
* THYROID 
* PARATHYROID GLAND
* ADRENAL GLAND
* PANCREAS
* OVARIES 
* TESTES

\textbf{TOPIC 3: APPLIED SCIENCE | DAY 2}


Applied Science is the branch of science that uses scientific ideas to solve real-life problems, create technology, and improve everyday life. It is also  the usage of scientific knowledge to solve real problems. A key feature of applied Science is Problem oriented-focus. 
Applied Science appears in many fields: Computer Science, Electric Engineering, Materials Science, Chemical Engineering, Materials Science, Biomedical Engineering, Aerospace Engineering, and Data Science.

\textbf{TOPIC 4: PSORIASIS | DAY 2}


Psoriasis is a chronic (long-lasting) skin condition where the body's immune system makes skin cells grow too fast. Because of this, thick, red, and scaly patches appear on the skin.


Key points:


It is not contagious — you cannot catch it from someone.


Patches often appear on the elbows, knees, scalp, back, or other parts of the body.


Causes include genetics, immune system problems, and triggers like stress, infection, cold weather, or skin injury.


Treatments help control symptoms, such as creams, light therapy, and oral or injectable medicines.

\textbf{TOPIC 5: EARTHQUAKE PREPAREDNESS | DAY 2}
 An earthquake is the shaking of the Earth's surface caused by the sudden release of energy in the Earth's crust. This usually happens when tectonic plates move or converge, diverge, and slide past each other creating stress that is released as vibrations called seismic waves.


Earthquake preparedness is the practice of getting ready before, during, and after an earthquake to stay safe and reduce danger. Before an earthquake, people should prepare an emergency kit with water, food, flashlights, medicines, and important documents. They must also secure heavy furniture and identify safe spots like under sturdy tables. During an earthquake, the safest action is to Drop, Cover, and Hold On, staying away from windows and heavy objects. After the shaking stops, it is important to check for injuries, avoid damaged buildings, turn off gas or electricity if needed, and stay alert for aftershocks. Being prepared helps save lives, prevents injuries, and reduces panic during disasters.

\textbf{TOPIC 1: THE CELL | DAY 3}


Cells are the basic building blocks of all living things.
Every plant, animal, and human is made up of cells. Some living things have only one cell (like bacteria), while others have millions or even trillions of cells (like humans).


What cells do:
* They give structure to living things.
* They take in nutrients and turn them into energy.
* They help living things grow and repair damaged parts.


* They reproduce to make new cells.




Two main types of cells:


* Animal cells
* Plant cells (these have extra features like a cell wall and chloroplasts)

\textbf{TOPIC 2: THE ENDOCRINE SYSTEM | DAY 3}


The endocrine system is made up of glands that produce hormones. These hormones travel through the blood and control many body activities such as growth, metabolism, mood, and reproduction. Major glands include the pituitary, thyroid, adrenal glands, pancreas, ovaries, and testes. This system helps keep the body balanced and functioning properly.

\textbf{TOPIC 3: APPLIED SCIENCE | DAY 3}


Applied Science is the use of scientific knowledge to solve real-world problems. It takes ideas from biology, chemistry, physics, and other sciences and applies them to create useful tools, technologies, and solutions. Examples include making medicines, designing machines, improving farming, and developing new materials.


Examples of Applied Science:


1. Medicine – Developing vaccines, antibiotics, and medical equipment.


2. Engineering – Designing bridges, cars, airplanes, and machines.


3. Agriculture – Creating fertilizers, pest control methods, and irrigation systems.


4. Environmental Science – Treating polluted water, recycling, and renewable energy solutions.


5. Technology – Building smartphones, computers, and household appliances.

\textbf{TOPIC 4: PSORIASIS | DAY 3}


Psoriasis is a chronic (long-term) skin disorder in which the skin cells grow too quickly, forming thick, red, scaly patches. Normally, skin cells grow and shed in about a month, but in psoriasis, this process happens in just a few days. This rapid buildup causes irritation, redness, and sometimes itching or pain. Psoriasis is not contagious, but it is linked to the immune system and genetics. Triggers can include stress, infections, cold weather, or certain medications. Treatment focuses on reducing inflammation, controlling symptoms, and slowing down skin cell growth through creams, ointments, light therapy, or medications.

\textbf{TOPIC 5: EARTHQUAKE PREPAREDNESS | DAY 3}


Earthquake preparedness means getting ready in advance to protect yourself, your family, and your property during an earthquake. It involves knowing what to do before, during, and after an earthquake to reduce injuries, save lives, and minimize damage.


Examples of earthquake preparedness:
* Making an emergency plan and practicing drills


* Preparing an emergency kit with food, water, and first-aid supplies


* Securing heavy furniture and appliances to prevent them from falling


* Knowing safe spots in your home (like under a table or against an interior wall)

\textbf{TOPIC 1: THE CELL | DAY 4}


The cell is the basic unit of life. It is the smallest structure that can perform all the functions needed for an organism to live, such as getting energy, growing, and reproducing. All living things—plants, animals, and even humans—are made up of cells. Some organisms have only one cell (unicellular), while others have many cells (multicellular).


Each cell has parts that help it work:
* The cell membrane acts like a protective gate.
* The cytoplasm is the jelly-like substance where activities happen.


* The nucleus controls the cell's actions.


* The mitochondria provide energy.


* In plants, chloroplasts make food, and the cell wall gives shape.


In short: The cell is the building block of life, and without cells, no living thing can exist.

\textbf{TOPIC 2: THE ENDOCRINE SYSTEM | DAY 4}
The endocrine system is a group of glands that release hormones into the bloodstream. These hormones act as chemical messengers that help control important body processes. The endocrine system regulates growth, metabolism, mood, stress responses, and reproductive functions. Some of the major glands include the pituitary gland, which directs other glands; the thyroid gland, which controls how fast the body uses energy; the adrenal glands, which help the body handle stress; the pancreas, which regulates blood sugar; and the ovaries or testes, which produce hormones for reproduction. By working together, these glands keep the body balanced and functioning smoothly.

\textbf{TOPIC 3: APPLIED SCIENCE | DAY 4}


Applied Science focuses on using scientific knowledge to solve real-world problems. It takes ideas from fields like biology, physics, and chemistry and applies them to create useful technologies, tools, and systems. Examples include medical innovations, engineering designs, environmental solutions, and everyday devices like smartphones and electricity. Applied Science helps improve daily life, health, safety, and technology through practical applications of scientific principles.

\textbf{TOPIC 4: PSORIASIS | DAY 4}


Psoriasis is a chronic (long-lasting) skin condition that causes the skin cells to grow too quickly, forming thick, red, scaly patches. It can appear on any part of the body but is most common on the elbows, knees, scalp, and back. Psoriasis is not contagious, but it can cause itching, pain, and discomfort. The exact cause is linked to the immune system and genetics. Treatments include creams, ointments, light therapy, and medications to reduce inflammation and control symptoms.

\textbf{TOPIC 5: EARTHQUAKE PREPAREDNESS | DAY 4}


Earthquake Preparedness is being ready in advance for an earthquake to keep yourself and others safe. It involves planning, practicing, and taking precautions so you can respond quickly when the ground shakes.


Examples of Preparedness Measures:


Emergency Plan: Know where to meet your family and how to communicate.


Safety Drills: Practice "Drop, Cover, and Hold On."


Emergency Kit: Keep water, food, flashlight, and first-aid supplies ready.


Home Safety: Secure furniture, appliances, and fragile items.


Awareness: Learn safe spots in your home, school, or workplace.

\textbf{TOPIC 1 : THE CELL | DAY 5}


The cell is the basic unit of life.
It is the smallest part of a living organism that can carry out all life processes such as growing, getting energy, and reproducing.


There are two main types of cells:
Plant cells – have a cell wall, chloroplasts, and a large vacuole.
Animal cells – have no cell wall and no chloroplasts.




Parts of a cell and their functions:
Cell membrane – protects the cell and controls what goes in and out.
Cytoplasm – jelly-like fluid where cell activities happen.
Nucleus – the "brain" of the cell; it controls everything.
Mitochondria – the "powerhouse"; gives energy to the cell.
Vacuole – storage area for water and nutrients.
Cell wall (plants only) – gives shape and support.
Chloroplast (plants only) – helps the plant make food through photosynthesis.

\textbf{TOPIC 2: THE ENDOCRINE SYSTEM | DAY 5}


The endocrine system is the body's chemical messaging system. It is made up of glands that produce hormones. These hormones travel through the blood and tell different parts of the body what to do.


What the endocrine system does:
* Controls growth and development
* Regulates metabolism (how the body uses energy)
* Maintains mood and emotions
* Controls reproduction
* Helps the body respond to stress


Main glands and what they do:
Pituitary gland – the "master gland"; controls other glands
Thyroid – controls body metabolism
Adrenal glands – help the body respond to stress
Pancreas – controls blood sugar
Ovaries/Testes – control reproduction and sexual development


The endocrine system keeps the body balanced and functioning by sending chemical signals called hormones.

\textbf{TOPIC 3: APPLIED SCIENCE | DAY 5}


Applied Science is the branch of science that uses scientific ideas, concepts, and discoveries to solve real-life problems. Instead of only studying how things work, Applied Science focuses on practical uses of science in everyday life.


It connects science to real-world applications such as:


1. Developing new medicines
2. Designing machines and technology
3. Improving agriculture
4. Creating safer buildings and transportation
5. Solving environmental problems


In short, Applied Science turns scientific knowledge into useful products, solutions, and technologies that help improve daily life.

\textbf{TOPIC 4: PSORIASIS | DAY 5}


Psoriasis is a long-term autoimmune skin disorder where the immune system speeds up the life cycle of skin cells. Instead of shedding normally, skin cells build up quickly, forming thick, red, scaly patches that may be itchy or painful. It is not contagious and often runs in families, meaning genetics play a role. Common triggers include stress, infections, certain medications, weather changes, and skin injuries.


Examples:
1. Plaque Psoriasis – The most common type, causing raised red patches with silvery scales on the elbows, knees, scalp, or back.


2. Guttate Psoriasis – Small red spots that appear after infections, like strep throat.


3. Nail Psoriasis – Causes pitting, thickening, or discoloration of fingernails or toenails.


4. Scalp Psoriasis – Red, scaly patches on the scalp that may extend beyond the hairline.

\textbf{TOPIC 5: EARTHQUAKE PREPAREDNESS | DAY 5}


Earthquake Preparedness means taking steps before, during, and after an earthquake to protect yourself, your family, and your property. Being prepared can save lives and reduce damage.


Key Steps for Preparedness:


Before an Earthquake:


* Make an emergency plan with your family.
* Prepare an emergency kit (food, water, first-aid supplies, flashlight).


* Secure heavy furniture and appliances.


* Know safe spots in your home


During an Earthquake:
* Drop, cover, and hold on.
* Stay indoors if safe; move away from windows and heavy objects.


* If outside, stay in an open area away from buildings.


After an Earthquake:
✓ Check for injuries and provide first aid.
✓ Be cautious of aftershocks.
✓ Inspect your home for damage and hazards.
}

\section{Statistical Analysis Summary}

\subsection{Comparative Assessment Validation}

{\color{maroon}
The dual-method validation between SimCSE and human teacher assessments yielded the following statistical measures:

\paragraph{SimCSE vs. Teacher Assessment:}
Mean absolute difference: 9.54 points (SD = 7.23)
Pearson correlation coefficient: r = 0.847
Statistical significance: p < 0.001
Cohen's d effect size: 1.24 (large effect)

\paragraph{Assessment Validity:}
SimCSE demonstrates exceptional alignment with teacher assessments, establishing strong validity for automated semantic similarity evaluation in educational contexts.
}

\section{System Architecture Specifications}

\subsection{Technical Implementation Details}

\paragraph{Frontend Technology Stack:}
React.js with Next.js framework for server-side rendering optimization, Tailwind CSS for responsive design implementation, Chart.js for real-time data visualization, and RESTful API integration for seamless data communication.

\paragraph{Backend Architecture:}
Node.js runtime environment with Express.js framework, PostgreSQL database with normalized schema design, SimCSE microservice integration through dedicated API endpoints, and comprehensive session management with secure authentication protocols.

\paragraph{Natural Language Processing Pipeline:}
BERT-base-uncased model serving as foundational encoder, SimCSE contrastive learning implementation with dropout-based data augmentation, cosine similarity computation for semantic distance measurement, and text preprocessing pipeline for formatting normalization.
}